\documentclass[12pt]{article}
\usepackage[utf8]{inputenc}
\usepackage[T1]{fontenc}
\usepackage{hyperref}
\title{Chol\'on Language Resources for K-12 Curriculum Development}
\author{Compiled by Claude Sonnet 5.5 (claude-sonnet-5-5) for Andr\'as Kornai}
\date{}
\begin{document}
\maketitle

\section{Introduction}
Chol\'on (ISO 639-3 cht) is a language of northern Peru, generally treated as extinct, though sources disagree about whether a rememberer survived past 2021. It is grouped with H\'ibito in the proposed Hibito--Cholon family. No school or revitalization materials were found.

\section{Primary Educational Resources}
None found.

\section{Linguistic and Reference Materials}
\begin{itemize}
\item Chol\'on language overview: \url{https://en.wikipedia.org/wiki/Chol%C3%B3n_language}
\item Hibito--Cholon languages overview: \url{https://en.wikipedia.org/wiki/Hibito%E2%80%93Cholon_languages}
\item M.\ Urban, article in \emph{Linguistic Discovery} on Chol\'on and the linguistic prehistory of northern Peru: \url{https://journals.dartmouth.edu/cgi-bin/WebObjects/Journals.woa/xmlpage/1/article/513}
\item Chapter in \emph{The Oxford Guide to the Languages of the Central Andes}: \url{https://academic.oup.com/book/59575/chapter/503209455}
\end{itemize}
Sources mention the 1748 \emph{Arte de la Lengua Cholona} by Pedro de la Mata and wordlists by Tessmann (1930); no URL for these was seen.

\section{Religious and Community Resources}
None found.

\section{Digital and Audio Resources}
None found.

\section{Curriculum Development Context}
Only historical and scholarly documentation exists. A heritage curriculum could at most be built around the colonial grammar and wordlists, with expert linguistic support.

\end{document}